\documentclass[10pt,a4paper]{amsart}
\usepackage{amsmath,amssymb,amsthm,graphicx}
\usepackage{verbatim}
\newenvironment{reference}{\comment}{\endcomment}
%\newenvironment{reference}{}{}
\newenvironment{slogan}{\comment}{\endcomment}
\newenvironment{history}{\comment}{\endcomment}

% For commutative diagrams we use Xy-pic
\usepackage[all]{xy}

% We use 2cell for 2-commutative diagrams.
\xyoption{2cell}
% Theorem environments.
%
\theoremstyle{plain}
\newtheorem{theorem}[subsection]{Theorem}
\newtheorem{proposition}[subsection]{Proposition}
\newtheorem{lemma}[subsection]{Lemma}

\theoremstyle{definition}
\newtheorem{definition}[subsection]{Definition}
\newtheorem{example}[subsection]{Example}
\newtheorem{exercise}[subsection]{Exercise}
\newtheorem{situation}[subsection]{Situation}

\theoremstyle{remark}
\newtheorem{remark}[subsection]{Remark}
\newtheorem{remarks}[subsection]{Remarks}

\numberwithin{equation}{subsection}

% Macros
%
\def\lim{\mathop{\mathrm{lim}}\nolimits}
\def\colim{\mathop{\mathrm{colim}}\nolimits}
\def\Spec{\mathop{\mathrm{Spec}}}
\def\Hom{\mathop{\mathrm{Hom}}\nolimits}
\def\Ext{\mathop{\mathrm{Ext}}\nolimits}
\def\SheafHom{\mathop{\mathcal{H}\!\mathit{om}}\nolimits}
\def\SheafExt{\mathop{\mathcal{E}\!\mathit{xt}}\nolimits}
\def\Sch{\mathit{Sch}}
\def\Mor{\mathop{\mathrm{Mor}}\nolimits}
\def\Ob{\mathop{\mathrm{Ob}}\nolimits}
\def\Sh{\mathop{\mathit{Sh}}\nolimits}
\def\NL{\mathop{N\!L}\nolimits}
\def\CH{\mathop{\mathrm{CH}}\nolimits}
\def\proetale{{pro\text{-}\acute{e}tale}}
\def\etale{{\acute{e}tale}}
\def\QCoh{\mathit{QCoh}}
\def\Ker{\mathop{\mathrm{Ker}}}
\def\Im{\mathop{\mathrm{Im}}}
\def\Coker{\mathop{\mathrm{Coker}}}
\def\Coim{\mathop{\mathrm{Coim}}}

% Boxtimes
%
\DeclareMathSymbol{\boxtimes}{\mathbin}{AMSa}{"02}

%
% Macros for moduli stacks/spaces
%
\def\QCohstack{\mathcal{QC}\!\mathit{oh}}
\def\Cohstack{\mathcal{C}\!\mathit{oh}}
\def\Spacesstack{\mathcal{S}\!\mathit{paces}}
\def\Quotfunctor{\mathrm{Quot}}
\def\Hilbfunctor{\mathrm{Hilb}}
\def\Curvesstack{\mathcal{C}\!\mathit{urves}}
\def\Polarizedstack{\mathcal{P}\!\mathit{olarized}}
\def\Complexesstack{\mathcal{C}\!\mathit{omplexes}}
% \Pic is the operator that assigns to X its picard group, usage \Pic(X)
% \Picardstack_{X/B} denotes the Picard stack of X over B
% \Picardfunctor_{X/B} denotes the Picard functor of X over B
\def\Pic{\mathop{\mathrm{Pic}}\nolimits}
\def\Picardstack{\mathcal{P}\!\mathit{ic}}
\def\Picardfunctor{\mathrm{Pic}}
\def\Deformationcategory{\mathcal{D}\!\mathit{ef}}
\usepackage{cleveref}
\crefname{originalremark}{original remark}{original remarks}
\Crefname{originalremark}{Original remark}{Original remarks}
\makeatletter
\newlabel{formal-spaces-remark-composition-variant-Noetherian}{{0GBK}{source}{}{sourcepointer-1}{}}
\newlabel{formal-spaces-remark-composition-variant-Noetherian@cref}{{[originalremark][0][]0GBK}{[1][0][]source}{}{}{}}
\makeatother
\crefname{originalreference}{original reference}{original references}
\Crefname{originalreference}{Original reference}{Original references}
\makeatletter
\newlabel{lemma-equivalent-with-artin-smooth}{{0GAJ}{source}{}{sourcepointer-2}{}}
\newlabel{lemma-equivalent-with-artin-smooth@cref}{{[originalreference][0][]0GAJ}{[1][0][]source}{}{}{}}
\makeatother
\crefname{originalreference}{original reference}{original references}
\Crefname{originalreference}{Original reference}{Original references}
\makeatletter
\newlabel{lemma-exact-sequence-NL}{{00S2}{source}{}{sourcepointer-3}{}}
\newlabel{lemma-exact-sequence-NL@cref}{{[originalreference][0][]00S2}{[1][0][]source}{}{}{}}
\makeatother
\crefname{originalreference}{original reference}{original references}
\Crefname{originalreference}{Original reference}{Original references}
\makeatletter
\newlabel{lemma-rig-smooth-continuous}{{0GCI}{source}{}{sourcepointer-4}{}}
\newlabel{lemma-rig-smooth-continuous@cref}{{[originalreference][0][]0GCI}{[1][0][]source}{}{}{}}
\makeatother
\crefname{originallemma}{original lemma}{original lemmas}
\Crefname{originallemma}{Original lemma}{Original lemmas}
\makeatletter
\newlabel{more-algebra-lemma-base-change-property-ext-1-annihilated}{{0G9H}{source}{}{sourcepointer-5}{}}
\newlabel{more-algebra-lemma-base-change-property-ext-1-annihilated@cref}{{[originallemma][0][]0G9H}{[1][0][]source}{}{}{}}
\makeatother
\crefname{originallemma}{original lemma}{original lemmas}
\Crefname{originallemma}{Original lemma}{Original lemmas}
\makeatletter
\newlabel{more-algebra-lemma-ext-1-annihilated}{{0G9K}{source}{}{sourcepointer-6}{}}
\newlabel{more-algebra-lemma-ext-1-annihilated@cref}{{[originallemma][0][]0G9K}{[1][0][]source}{}{}{}}
\makeatother
\crefname{originallemma}{original lemma}{original lemmas}
\Crefname{originallemma}{Original lemma}{Original lemmas}
\makeatletter
\newlabel{more-algebra-lemma-two-term-base-change}{{0G9G}{source}{}{sourcepointer-7}{}}
\newlabel{more-algebra-lemma-two-term-base-change@cref}{{[originallemma][0][]0G9G}{[1][0][]source}{}{}{}}
\makeatother
\title{Rig-smooth morphisms are stable under composition}
\author{The Stacks Project Authors}
\date{stacks-project:restricted.tex:local-snapshot-20261012; GNU FDL 1.2 or later}
\begin{document}
\maketitle
\noindent\textit{Editorial note.} All delivered mathematical text is the author's own, copied from the pinned Stacks Project chapter file restricted.tex; no statement or proof step was rewritten, reordered or removed.\par\medskip
\noindent\textit{Editorial note.} The counted proof is the author's complete argument: the proof is complete: the six-term sequence, the annihilation argument, the localization short exact sequence and the final criterion application are all written out\par\medskip
\noindent\textit{Editorial note.} The label lemma-composition-rig-smooth-continuous names the counted statement, and the proof environment immediately following it in the source is delivered verbatim as its proof.\par\medskip
\noindent\textit{Editorial note.} The counted statement is carried by the author's own lemma environment in the chapter source.\par\medskip
\noindent\textit{Allowed standard prerequisites (editorial).}\begin{itemize}
\item the rig-smooth criterion via H\textasciicircum{}0 of the completed naively trivialized cotangent complex being locally finite projective and H\textasciicircum{}\{-1\} killed by a power of the ideal of definition
\item the long exact cohomology sequence from the transitivity triangle for NL
\item the annihilator lemmas for the two-term complex under base change
\item localization and extensions of finite projective modules
\end{itemize}\medskip
\ifdefined\hypertarget\hypertarget{sourcepointer-1}{}\fi
\noindent\textit{Source reference.} Original remark 0GBK: sources/stacks-chow/tags.txt, source lines 20375--20375. Prerequisite pointer invoked at proof line(s) 4398; it names a standard fact of a sibling chapter that the counted argument uses, and every step of the counted proof built on it is retained verbatim in the delivered body.\par\medskip
\ifdefined\hypertarget\hypertarget{sourcepointer-2}{}\fi
\noindent\textit{Source reference.} Original reference 0GAJ: sources/stacks-restricted/pinned-original-source.tex, source lines 843--843. Prerequisite pointer invoked at proof line(s) 4423; it names a standard fact of the same chapter that the counted argument uses, and every step of the counted proof built on it is retained verbatim in the delivered body.\par\medskip
\ifdefined\hypertarget\hypertarget{sourcepointer-3}{}\fi
\noindent\textit{Source reference.} Original reference 00S2: sources/stacks-restricted/pinned-original-source.tex, source lines 725--725. Prerequisite pointer invoked at proof line(s) 4409; it names a standard fact of the same chapter that the counted argument uses, and every step of the counted proof built on it is retained verbatim in the delivered body.\par\medskip
\ifdefined\hypertarget\hypertarget{sourcepointer-4}{}\fi
\noindent\textit{Source reference.} Original reference 0GCI: sources/stacks-restricted/pinned-original-source.tex, source lines 4214--4214. Prerequisite pointer invoked at proof line(s) 4404; it names a standard fact of the same chapter that the counted argument uses, and every step of the counted proof built on it is retained verbatim in the delivered body.\par\medskip
\ifdefined\hypertarget\hypertarget{sourcepointer-5}{}\fi
\noindent\textit{Source reference.} Original lemma 0G9H: sources/stacks-chow/tags.txt, source lines 20302--20302. Prerequisite pointer invoked at proof line(s) 4426; it names a standard fact of a sibling chapter that the counted argument uses, and every step of the counted proof built on it is retained verbatim in the delivered body.\par\medskip
\ifdefined\hypertarget\hypertarget{sourcepointer-6}{}\fi
\noindent\textit{Source reference.} Original lemma 0G9K: sources/stacks-chow/tags.txt, source lines 20305--20305. Prerequisite pointer invoked at proof line(s) 4432; it names a standard fact of a sibling chapter that the counted argument uses, and every step of the counted proof built on it is retained verbatim in the delivered body.\par\medskip
\ifdefined\hypertarget\hypertarget{sourcepointer-7}{}\fi
\noindent\textit{Source reference.} Original lemma 0G9G: sources/stacks-chow/tags.txt, source lines 20301--20301. Prerequisite pointer invoked at proof line(s) 4425; it names a standard fact of a sibling chapter that the counted argument uses, and every step of the counted proof built on it is retained verbatim in the delivered body.\par\medskip
\begin{lemma}
\label{lemma-composition-rig-smooth-continuous}
The property $P(\varphi)=$``$\varphi$ is rig-smooth''
on arrows of $\textit{WAdm}^{Noeth}$ is stable under composition
as defined in Formal Spaces, Remark
\ref{formal-spaces-remark-composition-variant-Noetherian}.
\end{lemma}

\begin{proof}
We strongly urge the reader to find their own proof and not read the proof
that follows. The statement makes sense by
Lemma \ref{lemma-rig-smooth-continuous}.
To see that it is true assume we have rig-smooth morphisms
$A \to B$ and $B \to C$ in $\textit{WAdm}^{Noeth}$.
Then we can choose an ideal of definition $I \subset A$
such that the topology on $C$ and $B$ is the $I$-adic topology.
By Lemma \ref{lemma-exact-sequence-NL} we obtain an exact sequence
$$
\xymatrix{
C \otimes_B H^0(\NL_{B/A}^\wedge) \ar[r] &
H^0(\NL_{C/A}^\wedge) \ar[r] &
H^0(\NL_{C/B}^\wedge) \ar[r] & 0 \\
H^{-1}(\NL_{B/A}^\wedge \otimes_B C) \ar[r] &
H^{-1}(\NL_{C/A}^\wedge) \ar[r] &
H^{-1}(\NL_{C/B}^\wedge) \ar[llu]
}
$$
Observe that $H^{-1}(\NL_{B/A}^\wedge \otimes_B C)$
and $H^{-1}(\NL_{C/B}^\wedge)$ are annihilated by
a power of $I$; this follows from
Lemma \ref{lemma-equivalent-with-artin-smooth} part (2)
combined with
More on Algebra, Lemmas \ref{more-algebra-lemma-two-term-base-change} and
\ref{more-algebra-lemma-base-change-property-ext-1-annihilated}
(to deal with the base change by $B \to C$).
Hence $H^{-1}(\NL_{C/A}^\wedge)$ is annihilated by a power of $I$.
Next, by the characterization of rig-smooth algebras in
Lemma \ref{lemma-equivalent-with-artin-smooth} part (2)
which in turn refers to
More on Algebra, Lemma \ref{more-algebra-lemma-ext-1-annihilated} part (5)
we can choose $f_1, \ldots, f_s \in IB$ and $g_1, \ldots, g_t \in IC$
such that $V(f_1, \ldots, f_s) = V(IB)$ and
$V(g_1, \ldots, g_t) = V(IC)$ and such that
$H^0(\NL_{B/A}^\wedge)_{f_i}$ is a finite projective $B_{f_i}$-module and
$H^0(\NL_{C/B}^\wedge)_{g_j}$ is a finite projective $C_{g_j}$-module.
Since the cohomologies in degree $-1$ vanish upon localization at
$f_ig_j$ we get a short exact sequence
$$
0 \to
(C \otimes_B H^0(\NL_{B/A}^\wedge))_{f_ig_j} \to
H^0(\NL_{C/A}^\wedge)_{f_ig_j} \to
H^0(\NL_{C/B}^\wedge)_{f_ig_j} \to 0
$$
and we conclude that $H^0(\NL_{C/A}^\wedge)_{f_ig_j}$ is a finite
projective $C_{f_ig_j}$-module as an extension of same.
Thus by the criterion in
Lemma \ref{lemma-equivalent-with-artin-smooth} part (2)
and via that the criterion in
More on Algebra, Lemma \ref{more-algebra-lemma-ext-1-annihilated} part (4)
we conclude that $C$ is rig-smooth over $(A, I)$.
\end{proof}

\end{document}
